\section*{Chapitre \ref{ch:b3:radicals-carbenes} --- Radicaux, carbènes et réarrangements}

\begin{solution}{exo:b3:radicals-carbenes:1}
Le propane possède six hydrogènes primaires et deux secondaires. Chloration~: $6 \times 1 : 2 \times 3.9 = 6 : 7.8$,
soit 43\,\% de 1-chloropropane et 57\,\% de 2-chloropropane. Bromation~: $6 : 164$, 4\,\% et
96\,\%.
\end{solution}

\begin{solution}{exo:b3:radicals-carbenes:2}
$\ce{CH2}$~: triplet~; $\ce{CCl2}$~: singulet (les doublets non liants du chlore élèvent l'orbitale p vide).
Le dichlorocarbène s'additionne de façon stéréospécifique~: \textit{trans}-1,1-dichloro-2,3-diméthylcyclopropane.
\end{solution}

\begin{solution}{exo:b3:radicals-carbenes:3}
\textit{cis}-1,2-Diéthylcyclopropane~: le \omterm{def:b3:radicals-carbenes:carbenoid}{carbénoïde} s'additionne en une étape sur une seule face.
\end{solution}

\begin{solution}{exo:b3:radicals-carbenes:4}
Cyclohexanone~: $\varepsilon$-caprolactone (oxépan-2-one). Acétophénone~: acétate de phényle,
$\ce{CH3CO2C6H5}$~; le phényle migre de préférence au méthyle.
\end{solution}

\begin{solution}{exo:b3:radicals-carbenes:5}
Neuf hydrogènes primaires et un tertiaire. Chloration~: $9 : 5.1$, 64\,\% de 1-chloro-2-méthylpropane
et 36\,\% de 2-chloro-2-méthylpropane. Bromation~: $9 : 1600$, 0.6\,\% et 99.4\,\%.
\end{solution}

\begin{solution}{exo:b3:radicals-carbenes:6}
$k_c = 3.0 \times k_H[\mathrm{Bu_3SnH}] = 3.0 \times \num{2.4e6} \times 0.020 = \qty{1.4e5}{s^{-1}}$.
\end{solution}

\begin{solution}{exo:b3:radicals-carbenes:7}
Attaque sur C5~: 5-exo-trig, favorisée~; sur C6~: 6-endo-trig, favorisée
par les règles mais plus lente ici. 4-exo-tet~: favorisée. 5-endo-trig~:
défavorisée. 5-exo-dig~: favorisée. 3-exo-dig~: défavorisée.
5-endo-dig~: favorisée.
\end{solution}

\begin{solution}{exo:b3:radicals-carbenes:8}
La perte d'eau donne un cation tertiaire benzylique~; sur le carbone voisin,
un phényle et un méthyle pourraient migrer, et le phényle, d'aptitude plus
élevée, migre. Le cation stabilisé par l'oxygène perd un proton~:
3,3-diphénylbutan-2-one.
\end{solution}

\begin{solution}{exo:b3:radicals-carbenes:9}
Curtius~: $\ce{C6H5CON3 -> C6H5NCO + N2}$. Avec l'eau, l'acide carbamique perd $\ce{CO2}$~: aniline
(il se forme un peu de diphénylurée quand l'aniline rencontre de
l'\omterm{def:b3:radicals-carbenes:curtius}{isocyanate}). Avec l'éthanol~: \textit{N}-phénylcarbamate d'éthyle, $\ce{C6H5NHCO2C2H5}$.
\end{solution}

\begin{solution}{exo:b3:radicals-carbenes:10}
Chlore~: primaire $421.7 - 432 = \qty{-10}{kJ/mol}$, tertiaire $405.4 - 432 = \qty{-27}{kJ/mol}$~; brome~: $+56$ et
$+39$. Les deux arrachements par le chlore sont exothermiques, avec des
états de transition précoces qui ressentent peu la différence de
\qty{16}{kJ/mol}~; les deux arrachements par le brome sont endothermiques,
avec des états de transition tardifs dont les énergies d'activation
diffèrent de presque toute cette valeur~:
$\eu^{16000/(8.314 \times 298)} \approx 600$, l'ordre de grandeur du rapport observé.
\end{solution}

\begin{solution}{exo:b3:radicals-carbenes:11}
Oxime \textit{E} (OH en anti par rapport à C2, qui porte le groupe méthyle)~: C2 migre, l'azote s'insère
à côté de lui~: 7-méthylazépan-2-one. Oxime \textit{Z}~: C6 migre~: 3-méthylazépan-2-one. Dans le
\omterm{def:b3:radicals-carbenes:beckmann}{réarrangement de Beckmann}, c'est la géométrie de l'oxime qui décide~; dans
l'\omterm{def:b3:radicals-carbenes:baeyer}{oxydation de Baeyer--Villiger}, c'est l'\omterm{def:b3:radicals-carbenes:migratory}{aptitude migratoire}, et le carbone
secondaire migre~: 7-méthyloxépan-2-one.
\end{solution}

\begin{solution}{exo:b3:radicals-carbenes:12}
Fraction cyclisée $k_c/(k_c + k_H c)$~: 0.91, 0.49 et 0.09. Pour 95\,\%~: $k_Hc = k_c(1/0.95 - 1)$, donc
$c = 0.0526 \times \num{2.3e5}/\num{2.4e6} = \qty{5.0e-3}{mol/L}$. On la maintient en ajoutant l'hydrure lentement au
pousse-seringue, ou en utilisant une quantité catalytique d'halogénure
d'étain avec un réducteur qui régénère l'hydrure.
\end{solution}

\begin{solution}{pb:b3:radicals-carbenes:1}
\textbf{1.} $\ce{C6H10O + NH2OH -> C6H11NO + H2O}$.

\textbf{2.} L'acide est nécessaire pour éliminer sous forme d'eau le groupe
hydroxy de l'intermédiaire d'addition, mais trop d'acide protone
l'hydroxylamine, qui ne s'additionne plus.

\textbf{3.} Non~: les deux carbones liés au carbone de C=N sont équivalents
par la symétrie du cycle.

\textbf{4.} Dans l'oxime \textit{E}, le groupe OH pointe à l'opposé du carbone C2 porteur du méthyle (en anti
par rapport à lui)~; dans l'oxime \textit{Z}, vers lui.

\textbf{5.} L'inversion d'un azote d'oxime porteur d'un oxygène
électronégatif a une barrière élevée.

\textbf{6.} Il protone (ou sulfone) le groupe hydroxy, qui devient un bon
groupe partant, l'eau.

\textbf{7.} La liaison C--C du cycle en anti par rapport au groupe partant
migre du carbone vers l'azote, qui se trouve ainsi inséré dans le cycle~:
six atomes deviennent sept.

\textbf{8.} Un ion nitrilium cyclique.

\textbf{9.} $\ce{C6H11NO -> C6H11NO}$~: une isomérisation, de l'oxime au lactame.

\textbf{10.} Il s'agit du caprolactame, azépan-2-one, amide cyclique à sept
chaînons de l'acide 6-aminohexanoïque~; sa polymérisation par ouverture de
cycle donne le nylon-6.

\textbf{11.} \qty{98.0}{g/mol} et \qty{113.0}{g/mol}.

\textbf{12.} $\qty{1.0e6}{g}/\qty{98.0}{g/mol} = \qty{10204}{mol}$~; $\times 0.98 \times 0.98 = \qty{9800}{mol}$~; $\times \qty{113.0}{g/mol} = \qty{1.11e6}{g}$, \qty{1.11}{t}.

\textbf{13.} Oxime~: $\num{10204} \times 0.98 = \qty{1.00e4}{mol}$~; $\ce{H2SO4}$~: $\qty{1.30e4}{mol}$, ce qui donne autant de moles de
$\ce{(NH4)2SO4}$~: $\qty{1.30e4}{mol} \times \qty{132.1}{g/mol} = \qty{1.7}{t}$, plus que de caprolactame.

\textbf{14.} $\ce{C6H10O + RCO3H -> C6H10O2 + RCOOH}$.

\textbf{15.} L'adduit tétraédrique du peroxyacide sur le carbone du
carbonyle (l'intermédiaire de Criegee), un peroxyhémicétal.

\textbf{16.} Un groupe $\ce{CH2}$ du cycle (les deux sont équivalents) migre vers l'oxygène peroxydique
le plus proche~; le carboxylate $\ce{RCOO-}$ s'en va.

\textbf{17.} L'$\varepsilon$-caprolactone, oxépan-2-one~; sa polymérisation par ouverture de cycle donne la
poly($\varepsilon$-caprolactone), un polyester biodégradable.

\textbf{18.} Son seul sous-produit est l'eau, alors qu'un peroxyacide laisse
une mole d'acide carboxylique par mole de lactone et est lui-même dangereux
à stocker.

\textbf{19.} Le carbone secondaire C2 migre~: 7-méthyloxépan-2-one.

\textbf{20.} Oui~: le carbone qui migre conserve sa configuration.

\textbf{21.} 7-Méthylazépan-2-one (C2, en anti, migre).

\textbf{22.} 3-Méthylazépan-2-one (C6 migre).

\textbf{23.} Beckmann~: la géométrie de l'oxime (le groupe en anti migre).
Baeyer--Villiger~: l'\omterm{def:b3:radicals-carbenes:migratory}{aptitude migratoire} (le carbone le plus substitué
migre).

\textbf{24.} Ses deux carbones $\alpha$ sont équivalents~: une seule oxime, un seul lactame, une seule lactone.

\textbf{25.} Une tonne de cyclohexanone donne environ \qty{1.11}{t} de caprolactame.
\end{solution}
